% part: intuitionistic-logic
% chapter: propositions-as-types
% section: proof-terms

\documentclass[../../../include/open-logic-section]{subfiles}

\begin{document}

\olfileid{int}{pty}{ter}

\olsection{నిరూపణ పదాలు}

\Log{N1i}లోని ఊహలకు విడుదల గుర్తులుగా చరాలను (ఉదా., $!A^x$) కేటాయిస్తాం. \Log{N2i}లో ప్రతి సీక్వెంట్ $\Gamma \Sequent !A$కు ఎడమ సందర్భం~$\Gamma$లో చరం--\tetoken{సూత్రం}{!!{formula}} జతల జాబితా ఉంటుంది. అందువల్ల \Log{N2i}లో \tetoken{ఒక నిరూపణలోని}{!!a{proof}} ప్రతి సీక్వెంట్, ఇప్పటివరకు \tetoken{నిరూపణ}{!!{proof}} ఏమి \tetoken{నిరూపించిందో}{!!{prove}d} (అంటే $!A$) మాత్రమే కాక, ప్రతి దశలో ఏ ఊహలు తెరిచి ఉన్నాయో (అంటే $\Gamma$) కూడా తెలియజేస్తుంది. ప్రతి సీక్వెంట్‌లో $!A$ను \emph{ఎలా} \tetoken{నిరూపించారో}{!!{prove}d} కూడా చేరిస్తే ఎలా ఉంటుంది? ఇందుకోసం \emph{పదాలతో గుర్తించిన} వ్యవస్థ \Log{tN2i}ను నిర్వచించవచ్చు; దాని సీక్వెంట్ల రూపం $\Gamma \Sequent N : !A$. $\Gamma$, $!A$ మునుపటిలాగే ఉంటాయి: మొదటిది $x: !B$ రూపం గల జతల సమితి; రెండవది ఆ సీక్వెంట్‌తో ముగిసే ఉప-\tetoken{నిరూపణ}{!!{proof}}, సందర్భం~$\Gamma$లోని ఊహల నుంచి అనుసరిస్తుందని చూపిన \tetoken{సూత్రం}{!!{formula}}. $N$ అనేది ఆ సీక్వెంట్‌తో ముగిసే \tetoken{నిరూపణను}{!!{proof}} సంకేతీకరించే ఒక పదం.

ప్రతి సీక్వెంట్‌కు కేటాయించే పదాలను చరాలు $x$, $y$, \dots నుంచి, ఇప్పటివరకు చేసిన నిగమనాలను సంకేతీకరించే ప్రమేయ చిహ్నాలతో నిర్మిస్తాం. ప్రతి నిగమన నియమానికి వేరు ప్రమేయ చిహ్నం కావాలి. \tetoken{ప్రవేశపెట్టే}{!!{introduction}} నియమాలకు సంబంధించిన చిహ్నాలను “నిర్మాతలు”, \tetoken{తొలగింపు}{!!{elimination}} నియమాలకు సంబంధించినవాటిని “విచ్ఛేదకాలు” అంటాం. నియమానికి ఎన్ని ఆధారాలు ఉంటాయో చిహ్నానికి అన్ని ఆర్గ్యుమెంట్లు ఉంటాయి. ఉదాహరణకు \Intro{\land}, \Elim{\land}[1]లకు $c^\land$ (రెండు స్థానాలు), $d^\land_1$ (ఒక స్థానం) వాడవచ్చు. \Log{tN2i}లో ఆ నియమాలు:
\[
\Axiom$\Gamma_1 \fCenter N: !A$
\Axiom$\Gamma_2 \fCenter M: !B$
\RightLabel{\Intro{\land}}
\BinaryInf$\Gamma_1, \Gamma_2 \fCenter c^\land(N, M) : !A \land !B $
\DisplayProof
\quad
\Axiom$\Gamma \fCenter N: !A \land !B$
\RightLabel{\Elim{\land}[1]}
\UnaryInf$\Gamma \fCenter d^\land_1(N) : !A$
\DisplayProof
\]
నిగమన నియమం ఒక ఊహను విడుదల చేస్తే—అంటే ఆధార సందర్భంలో ఉన్న గుర్తించిన \tetoken{సూత్రం}{!!{formula}}~$x:!A$ నిర్ధారణ సందర్భంలో లేకపోతే—దీన్ని నిరూపణ పదంలో కూడా సూచించాలి. ఆ నియమాల ప్రమేయ చిహ్నాలు సంబంధిత చరాలను \emph{బంధిస్తాయి}. ఉదాహరణకు \Intro{\lif} నియమం $x:!A, \Gamma \Sequent !B$ నుంచి $\Gamma \Sequent !A \lif !B$కు తీసుకువెళుతుంది. నిర్మాణ చిహ్నం $c^{\lif}$ ఒక ఆర్గ్యుమెంట్‌ను తీసుకుంటుంది; కానీ చరం $x$ బద్ధమవుతుందని కూడా సూచించాలి. ఇలా నిరూపణ పదం, గుర్తు~$x$ గల ఊహ విడుదలైందని వ్యక్తపరుస్తుంది. ఆధార నిరూపణ పదం~$N$ అయితే, నిర్ధారణ పదాన్ని $c^{\lif}([x]N)$గా, నియమాన్ని ఇలా రాయవచ్చు:
\[
\Axiom$x:!A, \Gamma \fCenter N: !B$
\RightLabel{\Intro{\lif}}
\UnaryInf$\Gamma \fCenter c^{\lif}([x]N) : !A \lif !B$
\DisplayProof
\]
\sourcecorrection{OLTEPTPTYTER-001}{సంయోగ ప్రవేశ నియమం రెండవ ఆధార పదం $M$ అయినా, మూల నిర్ధారణలో చిన్న అక్షరం $m$ ఉంది; దాన్ని ఆధార పదంతో సమన్వయం చేశాం.}

\tetoken{ఒక నిరూపణను}{!!a{proof}} పూర్తిగా పునర్నిర్మించాలంటే అదనపు సమాచారం అవసరం. నియమం ఆధార \tetoken{సూత్రాలు}{!!{formula}s} దాని నిర్ధారణ \tetoken{సూత్రాన్ని}{!!{formula}} పూర్తిగా నిర్ణయించకపోతే, ఆ సమాచారాన్ని పదానికి చేర్చాలి. \Intro{\lif}లో ఇది అవసరం; అందువల్ల $c^{\lif}([x:!A]N)$ అని రాస్తాం. \Intro{\lor}, \FalseIntలలోనూ అవసరం కనుక ఆ సమాచారాన్ని మోసే ప్రమేయ చిహ్నాలను వాడతాం. ఉదాహరణకు:
\[
\Axiom$\Gamma \fCenter N:!A$
\RightLabel{\Intro{\lor}[1]}
\UnaryInf$\Gamma \fCenter c^{\lor{!B}}_1(N):!A \lor !B$
\DisplayProof
\quad
\Axiom$\Gamma \fCenter N:\lfalse$
\RightLabel{\FalseInt}
\UnaryInf$\Gamma \fCenter d^\lfalse_{!A}(N):!A$
\DisplayProof
\]
\sourcecorrection{OLTEPTPTYTER-002}{మూల వియోగ ప్రవేశ నిర్ధారణలో నిర్మాణ చిహ్నానికి ఆధార పదం $N$ ఆర్గ్యుమెంట్ లేదు; నిరూపణను సంకేతీకరించడానికి దాన్ని చేర్చాం.}

ఈ ఆలోచనలతో, ప్రతి సీక్వెంట్ రూపం $\Gamma
\Sequent N : !A$గా ఉండే సీక్వెంట్-శైలి సహజ నిగమన వ్యవస్థను నిర్మించవచ్చు; ఇక్కడ $N$ ఒక \emph{నిరూపణ పదం}.

ఇలాంటి నిరూపణ పదాలకు, \emph{రకాలతో కూడిన లాంబ్డా కలనశాస్త్రం}లోని పదాలకు సన్నిహిత సంబంధం ఉంది. ఆ సంబంధాన్ని గౌరవిస్తూ, పైనున్న సాధారణ నిర్మాణ/విచ్ఛేదక చిహ్నాలకు బదులు లాంబ్డా సాహిత్యంలో వాడే చిహ్నాలను తీసుకుంటాం. ఉదాహరణకు $c^{\lif}([x:!A]N)$కు బదులు $\lambd[\typeof{x}{!A}][N]$, $d^{\lif}(N,M)$కు బదులు $(NM)$, $c^\land(N,M)$కు బదులు $\pair{N}{M}$, $d^\land_i(N)$కు బదులు $\proj{i}{N}$ రాస్తాం.

\begin{defn}
  నిరూపణ పదాలను కింది నియమాలతో ఆగమన పద్ధతిలో ఉత్పత్తి చేస్తాం:
  \begin{enumerate}
  \item ప్రతి చరం $x$ నిరూపణ పదమే (అక్షయాలు).
  \item $x$ చరం, $!A$ \tetoken{ఒక సూత్రం}{!!a{formula}}, $N$ నిరూపణ పదం అయితే, $\lambd[\typeof{x}{!A}][N]$ కూడా నిరూపణ పదం~(\Intro\lif).
  \item $N$, $M$ నిరూపణ పదాలైతే, $(NM)$ కూడా నిరూపణ పదం~(\Elim\lif).
  \item $N$, $M$ నిరూపణ పదాలైతే, $\pair{N}{M}$ నిరూపణ పదం~(\Intro\land).
  \item $N$ నిరూపణ పదమైతే, $\proj{i}{N}$ కూడా నిరూపణ పదం; ఇక్కడ $i$ అనేది $1$ లేదా~$2$~(\Elim\land[i]).
  \item $N$ నిరూపణ పదం, $!A$ సూత్రం అయితే, $\inj[!A]{i}{N}$ నిరూపణ పదం; ఇక్కడ $i$ అనేది $1$ లేదా~$2$~(\Intro\lor[i]).
  \item $M$, $N_1$, $N_2$ నిరూపణ పదాలు, $x_1$, $x_2$ చరాలు అయితే, $\dcase{M}{x_1}{N_1}{x_2}{N_2}$ నిరూపణ పదం~(\Elim\lor).
  \item $N$ నిరూపణ పదం, $!A$ \tetoken{ఒక సూత్రం}{!!a{formula}} అయితే, $\abort{!A}{N}$ నిరూపణ పదం~(\FalseInt).
  \end{enumerate}
\end{defn}
\sourcecorrection{OLTEPTPTYTER-003}{వియోగ ప్రవేశ నిర్మాణంలో మూలం ప్రకటించిన నిరూపణ పదం $N$కు బదులు సంబంధం లేని $!M$ను వాడింది; దాన్ని $N$గా సరిచేశాం.}

ప్రతి నిరూపణ పదమూ \tetoken{ఒక నిరూపణకు}{!!a{proof}} సంబంధించిన పదం కాదు. ఉదాహరణకు $\pair{N}{M}$ పదం \Intro{\land} నిగమనంతో ముగిసే నిరూపణను సంకేతీకరిస్తుంది; అందువల్ల దాని చివరి-\tetoken{సూత్రం}{!!{formula}} సంయోగం~$!A \land !B$. అది \Elim{\lif} నియమానికి ప్రధాన ఆధారం కాలేదు కనుక $(\pair{N}{M}P)$ సరైన నిరూపణ పదం కాలేదు. \Log{tN2ip} నియమాలను పాటించే నిరూపణ పదాలే \emph{సరైన} నిరూపణ పదాలు. ఆ నియమాలు \olref{tab:tN2ip}లో ఉన్నాయి.

\subfile{rules-tN2}

\end{document}
